% 本文件由 LaTeX 排版 Agent 根据有效 translation.json 自动生成。
% author=Jerome; work_id=3005; title=驳斥路济弗尔派对话录

\chapter{驳斥路济弗尔派对话录}

% structure: p0001 source=h1 target=omitted reason=page-title-already-rendered-by-parent

这篇对话写于三七九年，即路济弗尔逝世后七年，且紧接热罗尼莫从加尔基斯旷野的隐修生活归来之后。虽然他接受了由保利诺（路济弗尔所祝圣者）的祝圣，但如本文所示，他对路济弗尔较狭隘的观点并无共鸣。路济弗尔原为撒丁岛卡利亚里的主教，约在三五四年因西方主教们试图谴责圣亚大纳削而首次崭露头角。他热忱地捍卫圣人的事业，并应其请求，由罗马主教利贝里乌斯派遣，偕同司铎潘克拉提乌斯与执事希拉里乌斯出使君士坦提乌斯皇帝。皇帝允准召开公会议，于三五四年在米兰举行。路济弗尔因抵制谴责亚大纳削的提议而声名大噪，且毫不畏惧地激烈反对皇帝。为此，他从三五五年至三六一年被流放，大部分时间居于巴勒斯坦的厄娄特洛波利，后转至忒拜德。正是此时，他的论战著作问世，其语气与格调仅凭书名即可见一斑：《\WorkTitle{\OriginalTerm{论背道君王}{De Regibus Apostaticis}}》与《\WorkTitle{\OriginalTerm{论不与异端者共融}{De non Conveniendo cum Hæreticis, etc.}}》。三六一年君士坦提乌斯驾崩后，儒利安准许流放主教回归；但路济弗尔并未前往亚历山大（那里将召开一次由亚大纳削主持的公会议，以弥合安提约基雅公教派内部的分裂——部分人追随默肋提，另一部分追随欧斯塔提），而是径直前往安提约基雅。在那里，他祝圣了后一派的首领保利诺为该教会的主教。味尔切的欧瑟伯随即携带亚历山大公会议的书信抵达，但见好事已先成，为避免与朋友冲突，便立即退去。路济弗尔留下，并\Quote{声明他拒绝与欧瑟伯或任何赞同亚历山大公会议温和政策的人共融。该公会议决定：实际的亚略派若放弃异端，可获宽恕，但不得担任教会职务；而仅同意亚略主义的主教则应予保留。正是后一项让步触怒了路济弗尔，从此他成为一项原则的捍卫者：凡曾向亚略主义作出任何妥协者，不得担任任何教会职务。}因此，他与亚大纳削本人也产生了对立——如前所述，亚大纳削主持了亚历山大公会议。最后，他返回撒丁岛的教区；据热罗尼莫编年史，他于三七一年去世。路济弗尔派在下一世纪初甚至更早便已消亡。它似乎从未形成一个独立的组织，尽管约三八四年有些路济弗尔派司铎曾向皇帝申诉，且盎博罗削与奥斯定都称他陷入了分裂。\par

本对话的论点可概述如下：如上所述，卡利亚里的路济弗尔曾在君士坦提乌斯统治时期因坚持亚大纳削的事业而被逐出教区；在儒利安即位宣布宽容后（三六一年），他前往安提约基雅，祝圣保利诺为主教。当时安提约基雅有三位主教：亚略派的多罗泰（三七六年继承尤佐依）、默肋提（虽信奉亚大纳削派见解，却由亚略派或半亚略派祝圣）以及保利诺；此外还有维大利，他是阿波林派团体的主教。路济弗尔出于反亚略派的激烈情绪，拒绝承认那些从亚略派归正的人为主教，但接受亚略派主教所施洗的平信徒。这一主张导致了路济弗尔派的分裂，并成为本对话的主题。\par

正统派在整个对话中指出：既然路济弗尔派承认亚略派主教所施行的圣洗有效，那么他们不承认那些悔改其亚略主张的主教便是前后矛盾的。路济弗尔派起初（2）急切地宣称亚略派无异于外邦人，但旋即意识到言过其实，于是收回此论。然而他声称（3）：接纳一个悔改的初学者是一回事，而接纳一个人为主教并举行圣体圣事则是另一回事。他说（4）：我们并不阻止堕落者悔改。正统派回答：我们通过接纳主教，不仅挽救了他们，也挽救了他们的羊群。路济弗尔派说（5）：\Quote{失了味的盐}不能再咸；又说：\Quote{基督与贝里雅耳有什么相通？}但对方回答（6）：这岂不证明了亚略派根本不能施行圣洗？是的，反对者说：他们如同洗者若翰，他的圣洗需要由基督的圣洗来补全。但对方回答：主教施行的是基督的圣洗，并赐予圣神。随后提到的坚振（9）与其说是必要的恩宠手段，不如说是教会的惯例。\par

争论似乎接近哲学上的文字游戏（10, 11），但路济弗尔派重拾论题。他说（12）：一个人单纯从亚略派主教接受圣洗，与理解异端的主教本人之间，有真正的区别。但对方回答（13）：两者一旦悔改，都应被接纳。\par

至此（14），路济弗尔派屈服了。但他希望确知正统派所推荐的确实是教会的实际做法。这引出了一段宝贵的教会史章节。正统派回顾了教会的胜利，而路济弗尔派却将这些胜利斥为败坏（15）。可耻的是，尽管他们持有真信经，但信德却太薄弱。接着他描述了（17,18）正统主教们如何被诱骗接受阿里米尼信经，但后来看到了自己的错误（19）。\Quote{世界呻吟地发现自己是亚略派。}他们尽一切可能纠正局面。为何不应接纳他们，就像在尼西亚接纳除异端发起人之外的所有人一样？（20）路济弗尔，那位好牧人，以及执事希拉里，将他们自己小小的团体分离成一个派别，却把其余的人留给了狼群（20,21）。麦子和稗子必须一起生长（22）。这始终是教会的原则（23），正如圣经（24）和宗徒传统所示，甚至西彼廉，当他希望悔改的异端者重新受洗时（25），也无法取胜。就连希拉里，通过从始终以补赎重新接纳异端者的教会领受圣洗（26,27），也承认了这一原则。我们有责任坚守在那教会及其分裂与实践中。\par

1. 不久前，一位路济弗尔的追随者与一位教会之子发生了争论。他的喋喋不休令人厌恶，所用语言极其恶毒。因为他宣称世界属于魔鬼，并如他们现今常说的那样，教会已变成了妓院。另一方则确实有理，但未顾及时间与地点，力主基督并没有白死，天主子从天上降下，绝非只为一件撒丁皮袍。简言之，争论尚未了结，黑夜打断了辩论，街灯亮起，那是集散会的信号。于是，斗士们退场，几乎要朝对方脸上吐唾沫；此前听众已商定在黎明时分于一处安静的廊柱下再会。他们果然都到了那里，并决定两位发言者的话应由记录员录下。\par

2. 众人就座后，路济弗尔派的赫拉弟乌说道：我要你先回答我的问题。亚略派是基督徒吗？\par

正统派：我用另一个问题回答你：所有的异端者都是基督徒吗？\par

路：你若称某人为异端者，就是否认他是基督徒。\par

正：那么，没有异端者是基督徒。\par

路：我先前已告诉过你。\par

正：他们若不属基督，就属魔鬼。\par

路：无人怀疑这一点。\par

正：但他们若属魔鬼，那么是异端者还是外邦人并无区别。\par

路：我不争论这一点。\par

正：那么我们一致同意，我们谈论异端者应如同谈论外邦人一样。\par

路：正是如此。\par

正：现在既然确定了异端者是外邦人，你随意提问吧。\par

路：我通过问题所欲引出的已明确陈述，即异端者不是基督徒。由此推论：若亚略派是异端者，而所有异端者都是外邦人，那么亚略派也是外邦人。但若亚略派是外邦人，且教会与亚略派（即与外邦人）毫无共融是毋庸置疑的，那么，你那接纳来自亚略派（即来自外邦人）的主教的教会，领受的便是卡皮托利山上的司祭而非主教，因此更应被称为反基督者的会堂而非基督的教会。\par

正：看！先知的话应验了：\BibleQuote{他们在我的面前挖了坑，自己反掉在其中。}\par

路：何以见得？\par

正：若亚略派如你所说，是外邦人，而亚略派的集会乃是魔鬼的营垒，那你为何接纳一个在魔鬼营垒中受洗的人？\par

路：我接纳他，但作为悔改者。\par

正：事实是你不知道自己在说什么。有谁接纳一个悔改的外邦人？\par

路：我出于单纯，开始时便回答所有异端者都是外邦人。但这个问题是刁钻的，你尽可在第一点上享有全胜。现在我进入第二点，主张：来自亚略派的平信徒，若悔改，应被接纳，但圣职人员则不可。\par

正：然而，你若承认我第一点，那么第二点也归我了。\par

路：请告诉我它如何归你。\par

正：难道你不知道圣职人员和教友只同有一位基督，且归化者与主教们并无不同的天主？那么，接纳平信徒的人为何不也接纳圣职人员？\par

路：为罪流泪与司掌基督身体有所不同；俯伏在弟兄脚前与在祭台上给子民分送圣体有所不同。为往事哀悼是一回事，在教会中放弃罪恶并活出光荣的生命是另一回事。你昨天还亵渎地宣称天主子是受造物，你每天都比犹太人更恶劣地以亵渎的石块投掷基督，你的双手沾满鲜血，你的笔是兵士的长矛；你，一时悔改的归化者，进入教会，如同奸夫进入童女一般？若你为自己的罪悔改，就放弃你的司铎职务；若你无耻于自己的罪，就保持原样。\par

正：你颇善辩，从论战的丛林飞向了演讲的开阔地。但我恳求你，避免陈词滥调，回到既定的范围和线索；之后，若你愿意，我们再扩大范围。\par

路：这并非演讲；我的愤慨难以忍受。随你如何陈述，随你如何辩论，你永远无法说服我：一个悔改的主教应被等同于悔改的平信徒。\par

正：既然你将一切归结为一点，并固执地坚持你的立场，即主教的情况与平信徒不同，我便依你所愿去做，且不会后悔利用你所提供的机会与你近距离交锋。请解释为何你接纳来自亚略派的平信徒，却不接纳主教。\par

路：我接纳承认自己曾堕入错误的平信徒；主不愿罪人死亡，而愿他悔改。\par

O. 那么也接纳一位主教，他与平信徒一样承认自己犯了错误，并且主不愿罪人死去，而愿他悔改，这句话依然有效。\par

L. 如果他承认自己的错误，为什么还要继续当主教？让他放下主教的职责，我就宽恕这悔改者。\par

O. 我要用你自己的话来回答你。如果一个平信徒承认自己的错误，他如何还能继续做平信徒？让他放下他的平信徒司祭职，也就是他的圣洗，我就宽恕这悔改者。因为经上记载：\BibleRef{默 1:6} \BibleQuote{祂使我们成为国度，成为祂的天主和父的司祭。} 又说：\BibleRef{伯前 2:9} \BibleQuote{圣洁的邦国，王家的司祭，特选的民族。} 一切禁止基督徒做的事，对主教和平信徒都是禁止的。做补赎的人谴责自己从前的生活。如果一个悔改的主教不能继续作他先前所是的，那么一个悔改的平信徒也不能留在他因之而承认自己是悔改者的那个状态中。\par

L. 我们接纳平信徒，因为没有人会被劝服改变，如果他知道他必须重新受洗。并且，如果他被拒绝，我们就成了他们丧亡的原因。\par

O. 接纳一个平信徒，你拯救一个灵魂；而我接纳一个主教，则使教会合一，我不说一座城的百姓，而是他为首领的整个省份；如果我赶走他，他会把许多人一同拖向丧亡。因此，我恳求你，将你认为有理由接纳几个人的同一理由，应用于全世界的得救。但是，如果你对此不满意，如果你如此苛刻，或者更确切地说，如此不讲道理地不仁慈，以致认为那施行圣洗的人是基督的仇敌，而你却把那领受圣洗的人当作儿子，我们并不如此自相矛盾：我们要么接纳主教，同时也接纳被他组成基督徒子民的百姓；要么我们不接纳主教，我们知道也必须拒绝他的百姓。\par

5. \Person{L}.：请问，难道你没有读过有关主教的话，\BibleRef{玛 5:13} \BibleQuote{你们是地上的盐；盐若失了味，可用什么使它再咸呢？它便毫无用处，只好丢在外面，任人践踏。}此外，还有司铎\BibleRef{肋 9:7}为有罪的百姓向天主转求，却无人为司铎代祷。现在这两处圣经都指向同一结论。因为正如盐调味一切食物，没有盐，什么食物也不能使人觉得可口；同样，主教是整个世界和他自己教会的调味品，如果他因否认真理、或异端、或情欲，或者一言以蔽之，因任何罪恶而失了味，那么当他是所有人的调味品时，还能用什么使他再咸呢？我们知道，司铎为平信徒献上自己的祭品，按手在顺服的他身上，呼求圣神的归来，这样，在邀请百姓的祈祷之后，把那个曾被交付于撒殚以毁灭肉体、使灵魂得救的人，重新与祭台和好；除非所有肢体都与他一同哭泣，他绝不恢复一个肢体的健康。因为当母亲为儿女恳求时，父亲容易宽恕儿子。那么，如果通过司铎的职务，一个忏悔的平信徒得以恢复与教会的关系，并且当忧伤在先时，宽恕随之而来，那么显然，一个被免去职位的司铎不能恢复他已丧失的地位，因为要么他成了忏悔者，就不能再是司铎，要么如果他继续任职，就不能通过忏悔纪律被带回教会。你竟敢用失了味的盐来败坏教会的味道吗？你竟敢把那被丢弃、本应躺在污泥里、被众人践踏的人放回祭台吗？那么宗徒的命令\BibleRef{铎 1:7} \BibleQuote{主教必须无可指摘，如同天主的管家}将会如何？再者，\BibleRef{格前 11:28} \BibleQuote{人应省察自己，然后才吃这饼、喝这杯}。我主的提示\BibleRef{玛 7:6} \BibleQuote{不要把圣物给狗，也不要把你们的珍珠投在猪前}又当如何？但如果你把这些话当作一般的劝诫，那么当对平信徒已经如此谨慎时，对司铎更当何等谨慎呢？\BibleRef{户 16:26} \BibleQuote{请你们离开这些恶人的帐幕，不要触摸他们的任何东西，免得你们因他们的一切罪恶而被消灭}，主藉梅瑟说。又在\WorkTitle{小先知书}中，\BibleRef{欧 9:4} \BibleQuote{他们的祭物必如居丧者的饼，凡吃这饼的必被玷污}。福音中主说，\BibleRef{玛 7:22} \BibleQuote{眼睛就是身体的灯；所以，如果你的眼睛纯洁，你的全身就充满光明。}因为当主教宣讲真信德时，黑暗就从所有人的心中散去。他给出了理由：\BibleRef{玛 5:15} \BibleQuote{人点灯，不放在斗底下，而是放在灯台上，就照亮所有在屋里的人。}也就是说，天主在主教心中点燃祂真知识的火，动机是使他不仅为自己发光，也为共同的益处发光。接着他说：\BibleRef{玛 6:23} \BibleQuote{如果你的眼睛有恶，你的全身就充满黑暗。所以，你里面的光如果黑暗了，那黑暗是多么大啊！}这是合理的；因为既然主教被设立在教会中是为了约束百姓远离错误，那么当教导者本人犯错时，百姓的错误将有多大？一个自身是罪人的人如何能赦免罪？一个不虔敬的人如何能使人成圣？当我的眼睛瞎了时，光如何能进入我内？啊，悲惨！假基督的门徒管理着基督的教会。对于这句话我们该想什么：\BibleRef{玛 6:23} \BibleQuote{没有人能事奉两个主人}？还有，\BibleRef{格后 6:14} \BibleQuote{光明与黑暗有什么相通？基督与贝里雅耳有什么相合？}在\WorkTitle{旧约}中我们读到，\BibleRef{肋 21:17} \BibleQuote{凡身上有缺陷的人，不可前来献上主的祭品。}又，\BibleQuote{亲近主自己的天主的司铎应洁净，恐怕主离弃他们。}又在同一处，\BibleQuote{当他们近前来在圣事中服务时，不可在自己身上带来罪恶，免得死亡。}还有许多其他段落，逐一举出将是无休止的，为简略起见，我略过它们。因为证明的效力不在于数量，而在于分量。这一切都证明，你们用一点面酵败坏了教会的整个面团，今天你们从昨天还像憎恶偶像一样憎恶的人手中领受圣体。\par

6. \Person{O}.：你的记忆很好，你确实以很大的篇幅给了我们许多圣经的引文：但绕遍了整个树林，你却被我的捕猎网捉住了。就算按你所说的，一个亚略派主教是基督的敌人，让他是失了味的盐，让他是没有火焰的灯，让他是失去瞳孔的眼睛：毫无疑问，你的论证会走到这一步——他自己没有盐，就不能调味别人；盲人不能照亮他人，也不能在自己光熄灭时点燃别人。但是，当你吞下他调味的食物时，为何责备那调味品无味？你的教会因他的火焰而明亮，你却控告他的灯熄灭了？他给你眼睛，而你却是瞎的？因此，我请你：要么既然你认可他的圣洗，就赋予他献祭的权力；要么如果你不认为他是司铎，就拒绝他的圣洗。因为不可能在圣洗中是圣洁的人却在祭台上是罪人。\par

L.：但是当我接受一个平信徒忏悔者时，是用按手和呼求圣神，因为我知道圣神不能被异端者给予。\par

正统派。你一切论点的途径都趋于同一交汇点，你就像受惊的鹿——当你逃避风中飘动的羽毛时，却陷入了最坚固的罗网。因为看到一个人在圣父、圣子和圣神的名下受洗，成为主的殿，而旧居被摧毁，为天主圣三建造了新圣所，你怎么能说在亚略人中间没有圣神的降临就可以赦免罪过呢？没有圣神的灵魂怎能洗去从前的污秽？因为洗濯灵魂的不仅仅是水，而是水本身先被圣神净化，才能以属灵的方式洗濯人的灵魂。\BibleRef{创 1:2} \BibleQuote{主的神}，\Person{梅瑟}说，\BibleQuote{运行在水面上}，由此可见没有圣神就没有圣洗。犹太地的\Place{贝特匝达}水池，若没有天使降临，就不能医治那些身体软弱者的四肢，而你竟敢把一个仅用普通水洗过的灵魂带到我面前，好像刚从浴池出来一样？我们的主\Person{耶稣基督}自己——更准确地说，不是祂被洗濯洁净，而是祂借自己的洗濯洁净了所有的水——祂的头刚从水中抬起，就领受了圣神。这不是说祂曾缺少圣神，因为祂是由圣神降孕而成为肉身的；而是为了证明真正的圣洗乃是圣神降临的圣洗。所以，如果一个亚略人不能赐予圣神，他甚至不能施洗，因为没有圣神就没有教会的圣洗。而你，当你接受一个由亚略人施洗的人，后来又呼求圣神，你要么应当给他施洗，因为若没有圣神他不能受洗，要么若他已在圣神内受洗，你就不可为你的皈依者呼求圣神，因为他在受洗时已经领受了祂。\par

7. 路济弗尔派。请问，你难道没有在\WorkTitle{宗徒大事录}中读到，那些已经由\Person{若翰}施洗的人，在回答宗徒的问题时说他们连圣神是什么都没有听说过，后来却领受了圣神？由此可见，受洗却可能没有圣神。\par

正统派：我不认为在场的听众对圣经如此无知，以致需要许多话来解决这个小问题。但在我说出任何支持我主张的话之前，请听我指出，按照你的观点，圣经中会引入怎样的混乱。我们说若翰在他的洗礼中不能将圣神赐予别人，却赐给了基督，这是什么意思？那位若翰是谁？\BibleQuote{在荒野中呼喊者的声音：预备主的道路，修直他的途径。}他曾说：\BibleRef{若 1:29}\BibleQuote{请看，天主的羔羊，除免世罪者！}我说得还不够：他从母腹中就喊道：\BibleRef{路 1:43}\BibleQuote{这是我哪里来的事，叫我主的母亲到我这里来？}他没有赐予圣神吗？而\BibleRef{宗 9:17}阿纳尼雅将圣神赐给了保禄吗？也许我把他置于众人之上似乎有些大胆。那么请听我主的话：\BibleRef{玛 11:11}\BibleQuote{凡妇人所生的，没有一个兴起比洗者若翰更大的。}因为没有一位先知既有福分宣布基督的来临，又用手指将他指出来。我有何必要详述这位如此卓越之人的赞美呢？当天主圣父甚至称他为天使时？\BibleRef{玛 11:10}\BibleQuote{看，我派遣我的使者在你面前，他要在你前面预备你的道路。}他必定是一位天使，因为在母腹中寄居后立即开始频繁出入荒野，还在婴儿时就与蛇嬉戏；当他一旦用眼睛注视基督后，便认为其他一切都不值得一看；他用配得天主使者的话音，宣讲主的言语，这些言语比蜂蜜和蜂房更甘甜。为不延迟我的问题，主的先驱理当如此成长。难道这样一位品性卓越的人不能赐予圣神，而百夫长科尔乃略却在受洗前就领受了圣神？请告诉我，他为什么不能赐予圣神？你不知道吗？那么请听圣经的教导：若翰的洗礼与其说在于得罪的赦免，不如说是一种悔改的洗礼，为得罪的赦免，即未来的赦免，这赦免将通过基督的圣化而实现。因为经上记载：\BibleRef{谷 1:4}\BibleQuote{若翰来了，在荒野中施洗，并宣讲悔改的洗礼，为使罪得赦免。}随后：\BibleRef{谷 1:5}\BibleQuote{他们在约但河里受他的洗，承认自己的罪。}因为正如他本人作为先驱在基督之前，他的洗礼也是主洗礼的前奏。\BibleRef{若 3:31}\BibleQuote{那从地上来的，属于地，所说的话也属于地；那从天上来的，超越万有。}又说：\BibleRef{玛 3:11}\BibleQuote{我固然以水洗你们，他却要以圣神洗你们。}但如果若翰，如他自己所承认的，没有以圣神施洗，那么他也没有赐予得罪的赦免，因为没有人能在没有圣神的情况下获得得罪的赦免。或者如果你争辩说，因为若翰的洗礼来自天上，所以得罪的赦免藉此得以赐予，那么请告诉我，我们在基督的洗礼中还能得到什么更多的东西？因它赐予得罪的赦免，它使人脱离地狱。因为它使人脱离地狱，它是完美的。但只有依赖于基督的十字架和复活的洗礼才能被称为完美的洗礼。尽管若翰自己说过：\BibleRef{若 3:30}\BibleQuote{他必须兴盛，我却必须衰微。}而你出于乖僻的挑剔，却给了仆人洗礼超过应得的，并毁了主人的洗礼，因为你没有给后者留下比前者更多的东西。你的主张有何用意？无非是——你不觉得奇怪吗，那些受若翰洗礼的人，后来藉着按手领受了圣神，尽管很明显他们甚至没有获得得罪的赦免，除非藉着将要来临的信德。但你却接受一个由亚略派人士施洗的人，并承认他有完美的洗礼，之后你又呼求圣神，仿佛这仍然是一个小小的缺陷，而实际上没有圣神的洗礼就不是基督的洗礼。但我走得太远了，当我本可以面对面地迎击对手、击退他的攻击时，却只是从远处投掷了几支轻矛。若翰的洗礼是如此不完全，以致显然那些受他洗礼的人后来接受了基督的洗礼。因为历史这样记载：\BibleQuote{阿颇罗在格林多的时候，保禄走遍了高原地带，来到厄弗所，遇见了一些门徒，问他们说：你们信从时，领受了圣神没有？他们回答说：我们连圣神是谁都没有听说过。保禄说：那么你们受的是什么洗？他们说：是若翰的洗礼。保禄说：若翰施的是悔改的洗礼，他告诉人民要信将来要来的那一位，就是要信耶稣。他们听了这话，就因主耶稣的名受了洗。保禄给他们按手，圣神便降在他们身上。}如果他们是因教会真实而合法的洗礼受洗，并因此领受了圣神，那么你就效法宗徒，为那些没有接受基督徒洗礼的人施洗，这样你也能呼求圣神了。\par

8. L. 口渴的人在梦中急切地痛饮溪流之水，越喝越渴。同样，在我看来，你似乎到处寻找论据来反驳我所提出的论点，却始终未能满足。难道你不知道圣洗后行覆手礼并呼求圣神，是教会的传统吗？你要圣经的证明吗？你可以在宗徒大事录中找到。即便它不以圣经的权威为基础，全世界的共识在这方面也具有命令的效力。因为教会的许多其他礼仪，源于传承，已获得了成文法般的权威，例如在洗礼池中三次浸头的做法，以及离开水后品尝搀和奶与蜜以象征婴孩的做法；又如主日站立敬拜、五旬节期间停止斋戒等习俗；还有许多其他未成文的做法，通过理性和习惯赢得了它们的地位。所以，你看我们遵行教会的惯例，尽管可能很清楚一个人是在呼求圣神之前受洗的。\par

9. O. 我不否认，对于那些远离大城镇、由长老和执事施洗的人，教会的惯例是主教前往探访，并通过覆手礼呼求圣神降临在他们身上。但我该如何形容你那种将教会法律应用于异端者，并将托付给你的贞女暴露在娼妓的妓院中的习惯呢？如果一位主教给人覆手，他是覆手在那些以正确信仰受洗、并相信父、子和圣神是三个位格却是一个本体的人身上。但亚略异端者只有这样的信仰（请听众掩耳，以免被如此粗鄙不敬的言语玷污）：唯独父是真天主，我们的救主耶稣基督是受造物，而圣神是二者的仆人。那么，一个尚未获得罪赦的人，如何能从教会领受圣神呢？因为圣神必须有一个洁净的住所；祂不会住进那个没有真信仰作为大司祭的殿宇。但如果你现在问，为何在教会中受洗的人，除非通过主教的手，否则不能领受我们宣称为真正圣洗中所赐予的圣神？让我告诉你，我们这一规则的权威来自于主升天后圣神降临于宗徒们的事实。在许多地方，我们发现这种做法更多是为了尊崇主教职，而非出于强制性的法律。否则，如果圣神只在主教的祈祷时才降临，那么那些在偏僻的房屋、堡垒或隐僻之地，由长老和执事施洗后，在主教授予前就已安眠的人，就太可怜了。一个教会的福祉取决于其大司祭的尊荣，若不赋予他某些特殊且独一无二的职能，那么教会中就会出现与司铎数量一样多的分裂。因此，没有圣职授予和主教的许可，长老和执事都没有施洗的权力。然而，若有需要，我们知道甚至平信徒也可以且经常施洗。因为人如何领受，也能如何给予；很难说我们必须相信\Person{斐理伯}施洗的太监缺少圣神。圣经这样论到他：\BibleQuote{两人下到水中，\Person{斐理伯}就给他施了洗。}离开水时，\BibleQuote{圣神降临在太监身上。}你也许认为应当以这段经文来反对，我们读到：\BibleQuote{那时，在耶路撒冷的宗徒们听说\Place{撒玛黎雅}接受了天主的圣言，就打发\Person{伯多禄}和\Person{若望}到他们那里去。他们二人一来到，就为他们祈祷，好使他们领受圣神，因为圣神还没有降临在他们任何人身上。}但为何如此，上下文告诉我们：\BibleQuote{他们只不过因主耶稣的名受了洗。于是宗徒们就给他们覆手，他们就领受了圣神。}如果你在这里说你做同样的事，因为异端者没有以圣神之名施洗，我必须提醒你，\Person{斐理伯}并没有与宗徒分离，而是属于同一教会，宣信同一位主耶稣基督；他无疑是后来为他的悔改者覆手的那些人的执事。但当你声称亚略异端者没有教会，只有会堂，他们的神职人员不敬拜天主，却敬拜受造物和偶像时，你怎能坚持认为在如此截然不同的情况下应当采用同一原则呢？\par

L. 你奋勇而坚定地正面抵挡了我的攻击；但你后方受创，背部暴露在箭矢之下。让我们甚至承认亚略异端者没有圣洗，因此圣神不能由他们赐予，因为他们自己尚未获得罪赦；这完全有利于我获胜，你所有的辩论挣扎不过是为我提供胜利的棕榈枝。亚略异端者没有圣洗，那么他们如何有主教职呢？他们中连平信徒都没有，怎么会有主教？我可能拒绝一个乞丐，你却要接受一个国王？你把营寨交给敌人，我们却要接纳他们中的一个逃兵吗？\par

11. 正：你若记得先前所说的话，就会知道自己已得到了答复；然而由于爱好反驳，你已偏离了主题，就像那些健谈而不善辩的人一样，无法争论时仍继续争吵。目前我的目标不是指控或辩护\OriginalTerm{亚略人}{Arians}，而是安全地绕过赛道的转折点，并坚持我们接纳一位主教，其理由与你们接纳一位平信徒相同。若你们赦免犯错者，我也宽恕悔改者。若按我们的信仰施洗的人对受洗者没有产生有害影响，那么按同一信仰祝圣主教的人也不会玷污被祝圣者。\OriginalTerm{异端}{Heresy}是微妙的，因此头脑简单的人容易被欺骗。被欺骗是平信徒和主教共有的命运。但你说，主教不可能犯错。事实上，被选为主教的人来自\Person{柏拉图}和\Person{阿里斯托芬}的怀抱。你能在他们中找到多少未完全受过这些作家教育的人？事实上，如今从识字阶层中被祝圣的所有人，无论他们是谁，其研究不是如何寻求圣经的精髓，而是如何用修辞的花朵来取悦民众的耳朵。此外还必须加上，\OriginalTerm{亚略异端}{Arian heresy}与世俗智慧携手并进，从\Person{亚里士多德}的源泉中借来论证的溪流。因此我们要像互相较量的孩童一样行事——无论你说什么，我都说；你断言什么，我就断言；你否认什么，我就否认。我们允许亚略人施洗；那么他必定是主教。若我们同意亚略人的洗礼无效，你们就必须拒绝平信徒，而我也不必接受主教。无论你走到哪里，我都跟随你；我们要么一同陷在泥潭里，要么一同脱身。\par

12. 路济斐：我们宽恕平信徒，因为当他受洗时，他真诚地以为自己是在加入教会。他相信了，并照其信仰受了洗。\par

正：一个人竟由非基督徒而成为基督徒，这倒是新鲜事。当他加入亚略人时，他按何种信仰受了洗？当然是按照亚略人所持的信仰。另一方面，若我们设想他自己的信仰是正确的，而他明知是由异端者施洗，那么他就不配得到我们给予犯错者的宽容。但若认为一个学生去向师傅学习，却在受教之前就懂得技艺，这是十分荒谬的。你能设想一个刚离弃偶像崇拜的人比他的老师更了解基督吗？若你说，他真诚地相信了圣父、圣子和圣神，因此获得了洗礼，那么请告诉我，真诚地对自己所信的无知是什么意思？他真诚地相信了。他相信了什么？当然，当他听到三个名号时，他相信了三个神，成了偶像崇拜者；或者因三个称号而相信一位有三个名字的神，从而陷入\OriginalTerm{撒贝略异端}{Sabellian heresy}。或者他可能受了亚略人的训练，相信只有一位真神，即圣父，而圣子和圣神是受造物。他还可能相信了什么，我不知道；因为我们很难认为一个在\Place{卡匹托尔}长大的人会学到\OriginalTerm{同体天主圣三}{co-essential Trinity}的教义。那样他就会知道圣父、圣子和圣神在本质上不可分，而在位格上可分。他也会知道子的名号包含在圣父的名号中，圣父的名号也包含在子的名号中。断言任何人在相信之前就能争论信仰，在入门之前就理解奥秘，施洗者和受洗者对天主的看法不同，这都是可笑的。此外，洗礼的习惯是，在宣认对三位一体的信仰之后，问：你相信圣教会吗？你相信罪过的赦免吗？你说他相信的是哪个教会？亚略人的教会？但他们没有教会。我们的教会？但此人并非受洗加入它；他不可能相信他所不知道的教会。\par

路济斐：我看得出来，关于我提出的每一点，你都能巧妙地喋喋不休；当我们射出一箭时，你便用一篇演说作为盾牌来避开；因此我将掷出一支长矛，它足以刺穿你的防御和你的连珠炮。我不会再让力量被诡计压倒。即使在教会之外受洗的平信徒，若他按信仰受洗，也只被接纳为悔改者；但主教要么不做补赎而继续做主教，要么若做补赎就不再是主教。因此我们既欢迎悔改的平信徒，又拒绝主教——若他希望继续其职务——这是正确的。\par

正：从拉紧的弓上射出的箭不易躲避，因为它在盾牌举起阻止之前便到达了目标。另一方面，你的命题没有锋芒，因此无法刺伤对手。那么你用尽全力掷出的、并说了如此威胁话语的那支矛，我按俗语所说的，用小指就把它拨开了。争议点不仅仅在于主教不能补赎而平信徒能够补赎，而是异端者是否领受了有效的洗礼。若他没有（根据你的立场必然如此），那么他在成为基督徒之前，如何能是悔改者？请向我证明，从亚略人那里来的平信徒有有效的洗礼，然后我才不否认他能够补赎。但若他不是基督徒，若他没有司铎使他成为基督徒，那么他尚未是信徒，如何能做补赎呢？\par

14. 路济斐：我恳求你放下哲学家的方法，让我们以基督徒的纯朴交谈；也就是说，若你愿意跟随的不是逻辑学家，而是加里肋亚的渔夫。你认为亚略人做主教合理吗？\par

正：你证明他是主教，因为你接纳他所施洗的人。而这就是你的过错之处：——既然我们在信仰上一致，并且都接纳亚略人，为什么我们之间还有隔阂的墙壁呢？\par

路济斐：我先前请求你不要像哲学家那样说话，而要像基督徒。\par

正：你是想学习，还是想争论？\par

路. 我当然辩论，因为我想知道你行为的理由。\par

正. 你若辩论，你已经得到了回答。我接纳一位亚略主教，正像你接纳一个只受过洗的人。你若愿意学习，请到我这边来：因为对手必须被战胜，只有门徒才能被教导。\par

路. 在我成为门徒之前，我必须听一位我觉得是我师傅的人讲道。\par

正. 你处理得不太公平：你希望我作你的老师，条件是你可以随时把我当作对手。因此，我要以同样的精神教导你。我们信仰一致，我们接纳异端者一致，让我们也在共融的条件上一致。\par

路. 那不是教导，而是辩论。\par

正. 就像你一手持盾要求和平，我也必须带着我的橄榄枝，其中嫁接了一把剑。\par

路. 我放下双手表示顺服。你是征服者。但在放下武器时，我问你强制我发的誓的意义。\par

正. 当然，但我首先恭喜你，并感谢基督我的天主，因你善良的心意使你从撒丁人的乏味教导转离，而接受那为全世界所公认的真理；不再像有些人那样说：\BibleQuote{主啊，救助；因为虔诚人消失了。} 他们以不敬之言使基督的十字架归于无效，将天主之子置于魔鬼的权下，并让我们如今将主对罪人的哀叹理解为适用于所有人：\BibleQuote{我流我的血有何益处呢？当我下到深渊时？} 但决不可让我们的主白白受死。\BibleRef{谷 3:27} 壮士已被捆绑，他的家产被夺走。圣父的话应验了：\BibleQuote{你向我请求，我必将万国赐你为基业，将地极赐你为田产。} \BibleQuote{于是水渠出现，大地的根基显露。} \BibleQuote{在其中他给太阳安设帐幕，没有一物可隐藏其热量。} 被天主完全充满的圣咏作者歌唱：\BibleQuote{敌人的刀剑已到尽头，你所倾覆的城邑也已毁灭。}\par

15. 我想知道，那些过分谨慎，或更确切地说，过分亵渎的人，他们断言会堂比教会更多，那么情况又如何呢？魔鬼的王国怎么被摧毁了，而在时代的终末，偶像怎么终于倾倒了？如果基督没有教会，或者如果他只有一座在撒丁岛的教会，那么他变得非常贫穷。如果撒殚同时拥有不列颠、高卢、东方、印度的各族、野蛮民族以及全世界，那么十字架的胜利品如何被收集在大地的一角呢？基督强大的对手，果然，把西班牙的蛇交给了他：他不屑于拥有一个贫穷的省份及其半饥半饱的居民。如果他们自诩福音的那节经文站在他们一边：\BibleRef{路 18:8} \BibleQuote{然而人子来的时候，遇得见世上有信德吗？} 让我提醒他们，这里所说的信德，正是主亲自所说的：\BibleRef{玛 9:22} \BibleQuote{你的信德救了你。} 以及别处关于百夫长的话：\BibleRef{玛 8:10} \BibleQuote{在以色列中，我从来没有遇见过这样大的信德。} 又对宗徒们说：\BibleRef{玛 8:26} \BibleQuote{为什么你们胆怯，你们这些小信德的人哪？} 在另一处也说：\BibleRef{玛 17:20} \BibleQuote{如果你们有信德像一粒芥菜子，你们对这座山说：从这边移到那边去，它就会移开。} 因为百夫长和那个患血漏十二年的穷妇人都没有相信天主圣三的奥迹，因为这些是在基督复活之后才启示给宗徒的；这样，相信天主圣三奥迹的人的信德才能有其应有的卓越；但主所赞许的是她单纯的心和对天主的虔诚：\BibleRef{玛 9:21} \BibleQuote{因为她心里说：只要我摸到他的衣边，我就必痊愈。} 这就是主说很少遇见的信德。这就是即使在那些信仰正确的人身上也很难完美找到的信德。\BibleRef{玛 9:29} \BibleQuote{照着你们的信德，给你们成就吧，} 天主说。说实话，我不喜欢那些话。因为如果照我的信德给我成就，我就必丧亡。然而，我确实相信天主圣父，相信天主圣子，相信天主圣神。我相信唯一的天主；尽管如此，我不愿照我的信德给我成就。因为仇敌常来，在主的麦田里撒下莠子。我并非暗示有什么比相信那奥迹的纯洁之心更大；但要找到对天主的无疑信德确实很难。为使我的意思清楚，我们假设一个情形：—— 我站着祈祷；如果我不相信，我就不能祈祷；但如果我真相信，我就该洁净我的心，因为那是看见天主的器皿；我就该用拳头捶打我的胸膛，眼泪会流下我的脸颊，我的身体会战栗，我的脸会变白，我应该伏在我主的脚前，用眼泪洗他的脚，用头发擦干，我应该紧抱十字架不放手，直到获得怜悯。但事实上，在祈祷中，我常常要么在回廊中行走，要么计算我的利息，要么被卑劣的念头带走，以致忙于一些我提起就脸红的事情。我们的信德在哪里？难道我们以为约纳是这样祈祷的吗？或者三个青年？或者在狮子坑中的达尼尔？或者在十字架上的强盗？我举这些例子是为了让你明白我的意思。但让每个人与自己的心交谈，他会在整个生命中发现在全然信德的灵魂是多么罕见，这样的灵魂不因爱慕虚荣而做任何事，不因人的闲言碎语而做任何事。因为禁食的人并不因此直接为天主禁食，向穷人伸手的人也不因此借贷给主。恶德紧邻美德。很难满足于只有天主作为审判者。\par

16. \Person{路济斐尔}说：我原本将那段经文留到最后用，你却抢先问了我。我们中几乎所有的人，更确切地说，已不再是我这一派的人，都把它当作辩论的攻城槌；看到它被击得粉碎，我非常高兴。但是，你能否好心地向我完整解释一下——不是以对手的身份，而是以门徒的身份——为什么教会接受来自亚略派的人？事实上，我无法回答你一个字，但我尚未从心里赞同你所说的话。\par

17. \Person{正统者}说：当\Person{君士坦提乌斯}在位，\Person{欧瑟比}和\Person{希帕提乌斯}担任执政官时，以合一与信仰为借口，制定了一个不忠的信经——如今已公认如此。因为在那时，似乎没有什么比追求合一、避免与世间其余的共融分离，更显虔诚，更合天主旨意。更何况当时的信仰宣认表面上已不再显示任何亵渎之词。他们说：\Quote{我们相信唯一真实的天主，全能圣父。我们也宣认：我们相信天主的独生子，祂在万世之前，在万有的起源之前，由天主所生。此外，我们相信独生子是单独由父单独所生，天主出自天主，相似于生祂的父，按照圣经；祂的诞生无人知晓，唯有生祂的父知道。}我们在其中能找到像\Quote{曾有一个时期，祂不存在}或\Quote{天主子是受造物，虽不是由存在之物所造}这样的话吗？没有。这无疑是完备的信德：说我们相信祂是天主出自天主。而且，他们称祂为独生子，\Quote{单独由父所生}。\Emph{生}是什么意思？当然是\Emph{非受造}。祂的出生消除了对祂是受造物的一切怀疑。他们进一步添加：\Quote{祂从天降下，因圣神受孕，由童贞玛利亚诞生，在般雀·比拉多手下被钉十字架，第三日从死者中复活，升天，坐在圣父的右边，祂将审判活人死人。}这些话中有虔诚的腔调，没有人想到毒药竟与如此宣告的蜜糖相混。\par

18. 关于\OriginalTerm{本体}{Usia}一词，它未被拒绝而没有理由。他们说：\Quote{因为它没有被发现于圣经中，}以及\Quote{它的新奇对许多人是一个绊脚石，我们认为最好不用它。}主教们并不在意这个名称，只要它所表达的内容得以保存。最后，正当谣言四起，说信仰声明中有虚假成分时，起草该声明的\Person{穆尔萨主教瓦伦斯}，在奉皇帝命令参加主教会议的禁卫军长官\Person{塔鲁斯}面前，宣称自己并非亚略派，而且极端厌恶他们的亵渎言论。然而，此事是秘密进行的，并未消除普遍的疑虑。因此，另一日，当许多主教和平信徒聚集在\Place{阿里米努姆}的教堂中时，\Place{比扎凯纳}省的主教\Person{穆佐尼乌斯}——因年长而被众人推为首位——发言如下：\Quote{我们中的一位已被授权向你们各位尊敬的教父宣读那些正在流传并已传到我们耳中的报告，以便那些本应刺耳并被逐出我们心中的恶毒意见能由我们同心一致地谴责。}全体主教回答说：同意。于是，当\Place{皮切努}省的主教\Person{克劳狄乌斯}应所有在场者之请，开始宣读归咎于\Person{瓦伦斯}的亵渎言论时，\Person{瓦伦斯}否认那是他的言论，并大声喊道：\Quote{若有人否认我们的主基督、天主子、在万世之前由父所生，让他受阿那特玛。}众人齐声赞同：\Quote{让他受阿那特玛。}\Quote{若有人否认圣子按照圣经相似于父，让他受阿那特玛。}全体回答：\Quote{让他受阿那特玛。}\Quote{若有人不说天主子与父同永，让他受阿那特玛。}众人再次齐声赞同：\Quote{让他受阿那特玛。}\Quote{若有人说天主子是受造物，如同其他受造物，让他受阿那特玛。}回答相同：\Quote{让他受阿那特玛。}\Quote{若有人说圣子是无中生有，却非出自天主父，让他受阿那特玛。}全体一同高喊：\Quote{让他受阿那特玛。}\Quote{若有人说：曾有一个时期圣子不存在，让他受阿那特玛。}这时，全体主教和整个教会一起以鼓掌和跺脚接受了\Person{瓦伦斯}的这番话。若有人以为我们编造了故事，让他查阅公开记录。无论如何，各教会档案柜中满是这些记录，此事在人们的记忆中犹新。一些参与此次主教会议的人仍然在世，连亚略派自己（这事实可以排除争议）也不否认我们叙述的准确性。因此，当众人把\Person{瓦伦斯}捧上天，并痛悔自己曾怀疑他时，先前开始宣读的\Person{克劳狄乌斯}说：\Quote{还有几点被我主及兄弟\Person{瓦伦斯}忽略了；若你们以为合适，为去除所有疑虑，让我们对它们通过一项普遍的谴责。若有人说天主子确实在万世之前，但绝不早于一切时间，从而将某物置于祂之前，让他受阿那特玛。}\Person{克劳狄乌斯}列举时，许多其他可疑之处也被\Person{瓦伦斯}谴责。若有人想了解更多详情，可在\Place{阿里米努姆}主教会议的记录中找到，这也是我本人的资料来源。\par

19. 这些程序结束后，大公会议解散。众人欢欣鼓舞地返回各自的教省。因为皇帝和所有善人都有一个共同的目标，即东西方应通过团结合一的纽带连接起来。但邪恶不会长久隐藏，在不良体液未清除之前就表面愈合的疮口会再次破裂。瓦伦斯、乌尔撒基乌斯以及与他们在邪恶中同谋的其他主教——他们当然曾是杰出的基督徒主教——开始挥舞他们的棕榈枝，并声称他们没有否认基督是受造物，而是说他与其他受造物相似。就在那一刻，\OriginalTerm{欧西亚}{\GreekText{οὐσία} (Usia)}一词被废除；尼西亚信仰被一致谴责。整个世界呻吟，并惊讶地发现自己成了亚流派。因此，有些人留在自己的共融中，另一些人开始给那些作为阿塔纳修斯追随者而流亡的宣信者写信；一些人绝望地哀叹他们曾进入的更好关系。但少数人，本于人性，将自己的错误辩护为智慧的展示。宗徒的船处于危险中，被风吹动，船舷被波浪拍打：再无希望。但主醒来，命令风暴平息；野兽死去，再次恢复平静。更明确地说，所有被逐出教座的主教，因新皇帝的仁慈，回到了他们的教会。随后埃及欢迎胜利的阿塔纳修斯；随后希拉里从战场返回高卢教会的怀抱；随后欧塞比乌斯返回，意大利脱下了哀悼的外衣。那些在里米尼陷入陷阱并无意中被报告为异端的主教开始聚集，同时他们以我们主的身体和教会中一切神圣事物为证，声称他们对自己的信仰没有任何错误的怀疑。他们说：\Quote{我们认为这些词应按其自然意义理解，我们从未怀疑在教会中——这真理宣认的纯朴与真诚之家——心里会隐藏一件事，嘴上说出另一件。我们对坏人评价过高，被欺骗了。我们没想到基督的主教们竟在反对基督。} 还有许多他们含泪说出的话，但为简洁起见，我略过不提。他们准备谴责之前的签署以及亚流派的所有亵渎言论。在此，我问我们那些过于苛求的朋友：他们认为对做出此宣认的人应如何处理？他们会说：\Quote{撤去老主教，祝圣新主教。} 这个方案被尝试过。但有多少良心不谴责自己的人会允许自己被撤职？尤其是当所有爱戴其主教的人们聚集起来，准备用石头砸死那些试图撤职他们主教的人。或许有人说，主教们本应留在自己的共融之内。也就是说，以愚蠢的残酷，他们将整个世界交于魔鬼。为何谴责那些并非亚流派的人？为何在教会在信德中保持和谐时撕裂教会？最后，难道他们要通过固执已见把正统信徒变成亚流派吗？我们知道，在针对亚流派背信而召集的尼西亚大公会议上，曾有八位亚流派主教被接纳，而当今世界没有一位主教的祝圣不是依赖于那次大公会议。既然如此，他们怎能反对它？——对它的忠诚曾使他们忍受流放之苦。\par

20. 路：\Quote{那么亚流派的人真的在之后被接纳了？请告诉我他们是谁。}\par

奥：\Quote{尼科米底亚主教欧塞比乌斯、尼西亚主教特奥格尼斯、当时为利比亚长老的萨拉斯、巴勒斯坦凯撒利亚主教欧塞比乌斯，以及其他不胜枚举的人；还有长老阿里乌斯——所有麻烦的最初源头；执事欧佐伊乌斯——他继欧多克西乌斯之后成为安提约基亚主教；以及读经员阿基拉斯。这三名亚历山大里亚教会的神职人员是异端的发起者。}\par

路：\Quote{假如有人否认他们曾被重新接纳，如何驳斥他？}\par

奥：\Quote{仍有曾参与那次大公会议的人健在。若这还不够——因为时间流逝，他们已很少，且证人不可能无处不在——若我们阅读尼西亚大公会议主教们的法令和名单，会发现我们刚才看到被重新接纳的那些人，确实与其他主教一同签署了\OriginalTerm{同质}{\GreekText{ὁμοούσιον} (homoousion)}。}\par

路：\Quote{你能指出尼西亚大公会议后，他们如何再次堕入不忠吗？}\par

正统派：\Quote{好建议，因为不信者惯于闭眼否认他们所不喜欢的事曾经发生。但他们之后除了重蹈覆辙还能做什么呢？既然大公会议是因他们而召集，而他们在会议前发表的书信和不敬的论著至今仍存留？因此，当那时三百多位主教欢迎了一些人——他们本可以拒绝这些人而不损害教会——我惊异某些人，他们确实维护尼西亚信德，却如此严厉，竟认为三位从流放归来的宣信者，为了世界的救恩，不应当做那么多杰出人物甘愿去做的事。但是，回到我们的出发点：宣信者归来后，在后来的\Place{亚历山大利亚}会议中决定，异端的发起者除外（他们不能以错误为由被宽恕），悔罪者应被接纳进入教会的共融：并非因为他们曾是异端者而能成为主教，而是因为那些被接纳的人显然并非异端者。西方同意了这个决定，正是通过这个因时代需要而达成的结论，世界才从撒殚的口中被夺回。我触及一个非常困难的主题，在此我不得不违背自己的意愿和目的，对那位圣德之人\Person{路济弗尔}有与他的功德及我的礼数所要求的不同看法。但我能做什么？真理开启我的口，催促我勉强的舌头说出心中的想法。在教会如此危急之时，当豺狼疯狂肆虐，他分开了少数羊只，却抛弃了羊群的余下部分。他本人是好牧人，但他却将大量战利品留给了猛兽。我不理会某些恶言者所传的报告，尽管他们坚称这些都是确凿的事实；例如，他这样做是出于对光荣的爱，以及留名于后世的渴望；或者又是，他因\Place{安提约基雅}的争执而对\Person{欧瑟比}心怀怨恨。对这样的人，我不相信任何这些报告；并且我现在仍将不断肯定这一点——我们与他之间的区别在于言语，而非实质，如果他真的接纳那些受亚略派施洗的人的话。}\par

21. 路济弗尔派：\Quote{我过去常听说的关于这些事的说法大相径庭，并且如我现在所想，更易于助长错误而非希望。但我感谢基督、我的天主，将真理之光倾注我心，使我不再不敬地称那属于祂的童贞教会为魔鬼的娼妓。还有一点我愿请你解释：关于\Person{希拉利}，他甚至连那些由亚略派施洗的人也不接纳，我们该说什么？}\par

正统派：\Quote{既然\Person{希拉利}离开教会时仅是一名执事，而且既然教会在他看来——尽管只在他个人看来——只是一群世俗的众人，他既不能适当地举行圣体圣事（因为他没有主教或司铎），也不能在没有圣体圣事的情况下施行圣洗。并且既然此人现已去世，因为他是一名执事，不能祝圣任何人继承他，他的派别也随他而亡。因为没有主教就不成其为教会。但略过那些极其少数、自以为是平信徒又是主教的人，让我向你指出我们对整个教会应持的看法。}\par

路济弗尔派：\Quote{你像俗语所说，用三个字解决了一个大问题。确实，当你说话时，我觉得我站在你一边。但当你停下来时，一些旧的疑虑又出现了：为什么我们接纳那些由异端者施洗的人？}\par

正统派：\Quote{这正是我所说的要指出我们对整个教会应持的看法时所想的。因为许多人都因你所说的疑虑而困惑。我的解释也许冗长，但如果真理得以彰显，那还是值得的。}\par

22. 诺厄的方舟是教会的预像，正如宗徒伯多禄所说——\BibleRef{伯前 3:20} \Quote{在诺厄的方舟内，只有少数，即八个人，藉着水得了救；这水所预表的圣洗，如今也拯救你们。}如同方舟内有各种走兽，在教会中也有各种种族和性格的人。如同方舟内有豹与山羊羔、狼与羔羊同处，同样在教会中也可找到义人和罪人，即\BibleRef{弟后 2:20} 金器银器与木器瓦器。方舟有舱房，教会也有许多住所。在诺厄方舟中八个人得救。而\BibleRef{训 11:2} \Quote{你要将一份分给七人，甚至八人，}即要相信两约。因此，有些圣咏标题写着\Emph{为第八日}，而第一百一十九篇被分为每段八节，每节以相应字母开头，用以训导义人。我们的主在山上对门徒所讲的八端真福，描绘了教会，也是八端。厄则克耳在建造圣殿时也用了这个数目。在圣经中你还会找到许多同样表达的事物。乌鸦从方舟中被放出，不再返回，随后鸽子向大地宣告和平。同样，在教会的圣洗中，那最不洁的鸟——魔鬼被驱逐出去，而圣神的鸽子向我们的大地宣告和平。方舟的结构是从三十肘宽开始，逐渐变窄至一肘。同样，教会由许多等级组成，最终归结为执事、长老和主教。方舟在洪水中受险，教会在这世上受险。诺厄离开方舟后栽种了葡萄园，饮了酒而醉。基督亦生于肉身，栽种了教会并受苦。大儿子嘲笑父亲的裸体，小儿子却遮盖了；犹太人嘲笑被钉的天主，外邦人却尊敬祂。若要我详论方舟的一切奥秘并与教会相比，白昼也不够用。我们中间谁是鹰？谁是鸽子和狮子？谁是鹿？谁是虫和蛇？就我们的主题所需，我将简要向你们指明。不仅绵羊住在教会中，也不仅洁净的鸟在那里飞翔；但在麦子中也有别的种子被撒下，\Quote{在整齐的麦田间，牛蒡、蒺藜和不结实的燕麦占据统治地位。}农夫该怎么办？拔掉稗子吗？那样连整个收成也一同毁掉了。农夫每天都用怪声驱赶飞鸟，或用稻草人吓唬它们；这里甩响鞭子，那里又摆出别的东西吓它们。然而，他仍受敏捷的狍子的侵袭或野驴的任性；这里老鼠把谷粒运到地下的仓库，那里蚂蚁密集地涌来践踏麦田。情况就是如此。没有一块田地是免于忧虑的。当家主睡觉时，仇敌在麦子中撒了莠子，仆人们提议去拔掉它们，主人却禁止了，他将簸扬谷糠和麦子的工作留给自己。在天主的家中有忿怒的器皿和怜悯的器皿，如宗徒所说。日子将到，教会的仓房被打开，主将取出忿怒的器皿；当他们离去时，圣人们将说：\BibleRef{若一 2:19} \Quote{他们从我们中间出去，却不是属于我们的；因为他们若是属于我们的，就必仍与我们同在。}没有人能自取基督的特权，在没有审判的日子之前，没有人能对人作出判决。如果教会已经洁净了，我们为主还保留什么呢？\BibleRef{箴 14:12} \Quote{有一条路，人以为正，至终却是死亡之路。}当我们的判断如此容易出错时，我们能依靠谁的意见呢？\par

23. 真福西彼廉曾试图避开破裂的水池，不饮异水；因此，他拒绝异端者的洗礼，召集了他的非洲主教会议，反对斯德望——斯德望是罗马宗座上有福的伯多禄的第二十二位继任者。他们聚集讨论此事，但尝试失败了。最后，那些曾与他一同决定必须重洗异端者的主教们，回归了古老的习惯，并颁布了一项新法令。你问我们必须遵循什么路线？我们所行的是我们先祖传给我们的，正如他们的先祖传给他们的一样。但何以要谈后来的时代？当基督的血刚刚流出，宗徒们仍在犹太时，主的身体被断言为幻影；迦拉达人被引去遵守法律，宗徒又一次为他们受产痛；格林多人不相信肉身的复活，宗徒以许多论据努力引他们归正。随后出现了西满术士和他的门徒默南得尔。他们自称是天主的德能。然后巴西理得发明了至高的神\OriginalTerm{阿卜拉克萨斯}{Abraxas}及其三百六十五种显现。然后尼苛劳，七位执事之一，他日夜不停纵欲，沉溺于污秽的幻梦。我不提基督来临前破坏传给他们法律的犹太异端者：多西托，撒玛黎雅人的首领，拒绝先知；撒杜塞人源自他的根，甚至否认肉身的复活；法利塞人因某些多余的仪式而自犹太人分离，从其异议的事实得名；黑落德党人接受黑落德为基督。我谈到那些撕裂了福音的异端者：撒图尔尼诺、敖斐特派、加因派、色特派、卡尔波克拉特、则林托及其继任者厄彼翁，以及其他的害虫，其中大多数在若望宗徒尚在时已爆发，然而我们不曾读到这些人中有谁被重洗。\par

24. 既然我们提及了那位杰出的圣人，也让我们从他的\WorkTitle{默示录}中表明，对异端者而言，没有圣洗的悔改应被视为有效。致厄弗所的天使的过失是他离弃了起初的爱心\BibleRef{默 2:4}。在培尔加摩教会天使的身上，吃祭偶像之物\BibleRef{默 2:14}和尼苛劳党人的道理\BibleRef{默 2:15}受到了谴责。同样，提雅提辣的天使因那自称是先知的依则贝耳、祭偶像之物和奸淫而受责斥\BibleRef{默 2:20}。然而，主鼓励所有这些人都要悔改，并且还加上了若不悔改将来必受惩罚的威胁。他若无意赦免悔改者，就不会催促他们悔改。难道有什么迹象表明他说过：\Quote{让那些在尼苛劳党人的信仰中受洗的人重领圣洗}？或者\Quote{让那些曾持有巴郎道理的培尔加摩人那时已信的人接受覆手}？不，反而他说：\Quote{所以你应当悔改\BibleRef{默 2:16}，不然我就要快来用我口中的剑攻击他们。}\par

25. 然而，如果那些由依拉略祝圣、近来成了无牧之羊的人，想要引证圣经来支持真福西彼廉在书信中留下的主张为异端者重洗的论点，我恳请他们记得：西彼廉并未诅咒那些拒绝跟随他的人。无论如何，他与反对他观点的人仍保持共融。他满足于劝诫他们，因诺瓦杜及当时许多其他异端的兴起，不要接纳任何不谴责其先前错误的人。事实上，他就此问题致罗马教宗斯德望的书信这样结论：\Quote{至亲爱的弟兄，因着我们彼此的尊重和真诚的爱，我将这些事告知你，我相信，出于你虔诚与信仰的真挚，你会赞同那些既合乎虔诚又真实的事情。但我知道有些人不愿放弃他们已持有的观点，且厌恶改变初衷；他们宁愿在不破坏同僚之间和平与和谐的纽带的情况下，一旦采纳了自己的计划就坚持不变。在这件事上，我们不强迫任何人，也不为任何人制定法律；让每个人在管理教会时跟随自己的自由选择：让每个人在自己的范围内做管理者，因为他必须向主交代自己的行为。}在致尤巴雅诺斯关于为异端者重洗的信末，他这样说：\Quote{我尽我微薄的能力写下了这些简短的话，并不想命令任何人，也不预判任何人的情况：我不愿阻碍任何一位主教在充分运用自己判断力的情况下做他认为正确的事。尽可能，我们避免与同僚和主教同侪就异端问题发生争辩，而是与他们保持神圣的和谐与主的平安，特别因为宗徒说：\BibleRef{格前 11:16}\InnerQuote{若有人似乎好争论，我们却没有这样的习惯，天主的各教会也没有。}我们以忍耐和温柔在心中保存爱德、我们品秩的尊荣、信德的纽带和主教职的和谐。}\par

26. 我还有另一个论证要提出，对此即使是现代的丢卡利翁依拉略也不敢喃喃一语。如果异端者未受洗因而必须重洗，因为他们原本不在教会内，那么依拉略自己也不是基督徒。因为他是在那个一直接纳异端洗礼的教会中受洗的。在阿利米农会议召开之前，在路济弗流亡之前，依拉略作为罗马教会的执事，曾接纳那些从异端归正的人，因他们先前已接受的洗礼。几乎没有可能只有亚略派是异端，而我们接纳所有除了他们施洗的人之外的人。你曾是一个执事，依拉略（教会可能这样说），接纳了摩尼派人施洗的人。你是一个执事，认可了厄比雍的洗礼。亚略兴起之后，你忽然对自己大失所望。你和你的家人与我们分离，并开设了你们自己的新洗池。如果有天使或宗徒重洗了你，我不会贬低你的做法。但既然你向我举剑，你是我的骨中之肉、我所哺育的，请归还我所给你的，并且如果你能，以其他方式做一个基督徒。假设我是一名妓女，我仍是你的母亲。你说我不守婚姻的床榻清白；但我现在是什么，在你受孕时我就是什么。若我与亚略犯奸淫，我之前同样与普拉克西亚、厄比雍、塞林托和诺瓦杜犯过奸淫。你敬重他们，欢迎他们，虽然他们是奸夫，到母亲的家中来。我不知道为什么一个奸夫比其他奸夫更令你生气。\par

27. 但若有人怀疑我们的祖先是否始终接纳异端者，让他读一读真福西彼廉的书信，他在信中责备了罗马主教斯德望及其因习用而根深蒂固的错误。也让他读一读依拉略为反对我们而发表的关于为异端者重洗的小册子，他会在那里发现依拉略自己承认：尤利、玛尔谷、西尔维斯特及其他古代主教同样接纳所有异端者悔改；而且，为了表明他不能正当地主张拥有真正的传统，不久前我们提到的尼西亚大公会议也接纳了所有异端者，除了撒摩沙塔的保禄的门徒。此外，它还允许一位诺瓦杜主教在归正后享有司铎的品秩；这一决定同时谴责了路济弗和依拉略，因为同一个人既被祝圣也被施洗。\par

28. 我也许可以整天讲同样的话，并以教会唯一的太阳\OriginalTerm{太阳}{Sun}枯竭一切论据的溪流。但既然我们已经讨论了很久，漫长的争论使听众的注意力疲倦了，我就无保留地简要告诉你我的意见。我们应当留在那由宗徒建立并延续至今的教会中。如果你曾听说有些自称为基督徒的人，他们的名称不是来自主耶稣基督，而是来自其他，例如，\Person{马尔西翁}派、\Person{华伦提努}派、山上或平原派，你可以肯定，那里不是基督的教会，而是反基督的会堂。因为他们是在教会建立之后才兴起的，这证明他们就是宗徒所预言要来的那些人。如果他们以为自己的主张有圣经的权威，就不要再自欺了，因为魔鬼自己也曾引用圣经，而圣经的精髓不在字句，而在意义。否则，如果我们遵循字句，我们也可以炮制一个新教条，主张穿鞋和有两件外衣的人不应被接入教会。\par

L. 你不要以为胜利只属于你。我们两人都是得胜者，每人赢得棕榈枝——你胜了我，而我胜了我的错误。但愿我每次辩论时，都能如此幸运地以更正确的意见取代错误的意见。不过，我必须承认，因为我最了解我一方的性格，他们更容易被征服，而不是被说服。\par

\SourceRecord{Translated by W.H. Fremantle, G. Lewis and W.G. Martley.；Nicene and Post-Nicene Fathers, Second Series；Vol. 6.；Edited by Philip Schaff and Henry Wace.；Buffalo, NY: Christian Literature Publishing Co.,；1893.；<http://www.newadvent.org/fathers/3005.htm>.}
